\chapter{عدة الحامل، وأحكام الإيلاء، وميراث الجد (ختام الكتاب)}

\section{عدة الحامل المتوفى عنها زوجها}
الحمد لله رب العالمين، والصلاة والسلام على نبينا محمد وعلى آله وصحبه أجمعين. أما بعد؛ فهذا هو الدرس الختامي من دروس شرح كتاب «الرسالة» للإمام أبي عبد الله الشافعي رحمه الله تعالى.

بدأ الشافعي رحمه الله بذكر مسألة اختلف فيها الصحابة قديماً، ليوضح كيف حسمت السنة هذا الخلاف. قال الله عز وجل: ﴿وَالْمُطَلَّقَاتُ يَتَرَبَّصْنَ بِأَنفُسِهِنَّ ثَلَاثَةَ قُرُوءٍ﴾، وقال: ﴿وَاللَّائِي يَئِسْنَ مِنَ الْمَحِيضِ مِن نِّسَائِكُمْ إِنِ ارْتَبْتُمْ فَعِدَّتُهُنَّ ثَلَاثَةُ أَشْهُرٍ وَاللَّائِي لَمْ يَحِضْنَ ۚ وَأُولَاتُ الْأَحْمَالِ أَجَلُهُنَّ أَن يَضَعْنَ حَمْلَهُنَّ﴾، وقال: ﴿وَالَّذِينَ يُتَوَفَّوْنَ مِنكُمْ وَيَذَرُونَ أَزْوَاجًا يَتَرَبَّصْنَ بِأَنفُسِهِنَّ أَرْبَعَةَ أَشْهُرٍ وَعَشْرًا﴾.

اختلف بعض أصحاب النبي ﷺ في الحامل المتوفى عنها زوجها؛ هل تعتد بأبعد الأجلين (أي تنتظر حتى تضع حملها وتمر عليها أربعة أشهر وعشراً)؟ أم تنقضي عدتها بوضع الحمل فقط؟ 
ظاهر القرآن يحتمل المعنيين، فدلت سنة رسول الله ﷺ على أن وضع الحمل هو آخر العدة مطلقاً. فقد روى الزهري أن سُبيعة الأسلمية رضي الله عنها وضعت حملها بعد وفاة زوجها بليالٍ قليلة، فمر بها أبو السنابل بن بعكك رضي الله عنه ورآها قد تصنعت للخطاب، فأنكر عليها وقال: إنها تعتد أربعة أشهر وعشرا. فذكرت ذلك لرسول الله ﷺ، فقال: «كَذَبَ أبو السنابل، قد حللتِ فتزوجي». 
(ومعنى "كذب" في لغة أهل الحجاز: أخطأ، أي أخبر بخلاف الواقع). فدلت السنة على أن وضع الحمل ينهي العدة فوراً.

\section{تقديم السنة على أقوال الرجال}
يستنبط الشافعي قاعدة عظيمة من هذه المسألة: "ما دلت عليه السنة فلا حجة في أحد خالف قوله السنة، ولو كان من أصحاب رسول الله ﷺ". فلا يحل لمسلم أن يقدم قول أحد من الناس على كلام الله وكلام رسوله ﷺ، أو على إجماع سلف الأمة. 

\section{أحكام الإيلاء (اليمين على ترك الجماع)}
ثم انتقل الشافعي لمسألة ليس فيها نص سنة مباشر، وإنما يُستدل عليها من القرآن واستنباطاً، وهي مسألة "المُولي" (الذي يحلف ألا يجامع امرأته).
قال تعالى: ﴿لِّلَّذِينَ يُؤْلُونَ مِن نِّسَائِهِمْ تَرَبُّصُ أَرْبَعَةِ أَشْهُرٍ ۖ فَإِن فَاءُوا فَإِنَّ اللَّهَ غَفُورٌ رَّحِيمٌ * وَإِنْ عَزَمُوا الطَّلَاقَ فَإِنَّ اللَّهَ سَمِيعٌ عَلِيمٌ﴾.

اختلف الصحابة؛ هل بمجرد انقضاء الأشهر الأربعة تطلق المرأة تلقائياً؟ أم يوقف الزوج ويُخيّر؟
رجح الشافعي أنه يوقف، فلا يلزمه الطلاق تلقائياً، بل إذا مضت الأشهر الأربعة قيل له: إما أن تفيء (تجامع) وإما أن تطلق. لأن الله خيره بين الفيئة والطلاق بعد مضي المدة، فلو كانت تطلق بمجرد انقضاء المدة لما كان هناك معنى للتخيير بعدها. والفيئة تكون بالجماع لمن قدر عليه، وهو يمين جعل الله له مدة لئلا تُضار الزوجة، فإن مضت أُجبر الزوج على أحد الأمرين.

\section{مسائل في المواريث (رد الفضل، وميراث الجد)}

\subsection{رد فضل المواريث}
اختلفوا فيما لو مات إنسان وترك أختاً، فلها النصف. فماذا يُفعل بالنصف الآخر؟
ذهب زيد بن ثابت رضي الله عنه والشافعي إلى أن النصف الباقي يذهب لجماعة المسلمين (بيت المال)، ولا يُرد عليها؛ لأن الله فرض لها النصف منفردة أو مع أخواتها، فلو رددنا عليها النصف الباقي لكنا ورثناها غير ما فرض الله لها. وذهب آخرون إلى الرد على ذوي الأرحام لقوله تعالى: ﴿وَأُولُو الْأَرْحَامِ بَعْضُهُمْ أَوْلَىٰ بِبَعْضٍ فِي كِتَابِ اللَّهِ﴾.

\subsection{ميراث الجد مع الإخوة}
اختلف الصحابة في الجد مع الإخوة؛ فذهب أبو بكر وابن عباس رضي الله عنهم إلى أن الجد كالأب يحجب الإخوة تماماً ويأخذ الميراث. وذهب زيد بن ثابت وعمر وعلي وابن مسعود رضي الله عنهم إلى أن الجد يُقاسم الإخوة ولا يحجبهم (بشروط وتفصيلات معروفة في الفرائض).
ورجح الشافعي مذهب المشاركة؛ وناظر من قال بالحجب بأن الجد يدلي بقرابة غيره (أبو أب الميت) والأخ يدلي بقرابة غيره (ابن أب الميت)، وكلاهما يدلي بقرابة الأب، فليس أحدهما أولى بالحجب من الآخر، فكان التشريك أوفر بالقياس وأقرب للعدل.

\section{ختام كتاب الرسالة}
ختم الإمام الشافعي كتابه بالتأكيد على أن العلم مراتب، وأن الاجتهاد يكون بالدليل، وأن التوفيق بيد الله.
وقد ختم الربيع بن سليمان (تلميذ الشافعي وراوية كتابه) نسخته قائلاً: "أجاز الربيع بن سليمان صاحب الشافعي نسخ كتاب الرسالة... في ذي القعدة سنة خمس وستين ومائتين وكتب الربيع بخطه".

ونحمد الله العلي القدير أن يسر لنا قراءة ومدارسة هذا الكتاب العظيم للإمام الشافعي طيب الله ثراه، ونسأله بأسمائه الحسنى وصفاته العلى أن يرزقنا العلم النافع والعمل الصالح، وأن يقسم لنا من خشيته ما يحول به بيننا وبين معاصيه. وصلى الله وسلم وبارك على نبينا محمد وعلى آله وصحبه أجمعين.